% ============================================================
% TERMINI DEL GLOSSARIO
% ============================================================


\section{Termini}

% -------------------- A --------------------

\phantomsection
\addcontentsline{toc}{subsection}{A}
\subsection*{A}

\begin{itemize}
    \item \textbf{Agile:} Approccio allo sviluppo software basato su iterazioni brevi, collaborazione continua e adattamento ai cambiamenti dei requisiti.
    \item \textbf{AIRA:} AI Requests Automation. Il nome identifica il progetto, che prevede la realizzazione di uno strumento multimodale basato sull'Intelligenza Artificiale (Artificial Intelligence, AI), in grado di acquisire ed elaborare informazioni provenienti da diversi canali di comunicazione, come e-mail, messaggi, note vocali e immagini. Il sistema interpreta i contenuti ricevuti, estrae le informazioni rilevanti e le converte in dati strutturati, integrandoli automaticamente con i sistemi gestionali aziendali.    \item \textbf{Analisi dei Requisiti:} Attività finalizzata a individuare, descrivere e documentare le funzionalità e i vincoli che il sistema deve soddisfare.
    \item \textbf{Analisi dei Rischi:} Attività che identifica i possibili eventi negativi per il progetto, ne valuta probabilità e impatto e definisce strategie di mitigazione.
    \item \textbf{API:} Application Programming Interface. Insieme di regole e strumenti che consente a componenti software differenti di comunicare e scambiarsi informazioni    
    \item \textbf{API REST:} Application Programming Interface basata su REST (Representational State Transfer). È un'interfaccia che permette a sistemi software differenti di comunicare attraverso richieste HTTP (Hypertext Transfer Protocol), utilizzando risorse identificate da specifici URL (Uniform Resource Locator).    \item \textbf{Attore:} Entità esterna al sistema, come una persona o un altro software, che interagisce con esso per raggiungere un determinato obiettivo.
    \item \textbf{Artefatto:} Prodotto concreto generato durante lo svolgimento del progetto, come un documento, un diagramma, un componente software o un file di configurazione.
\end{itemize}
% -------------------- B --------------------

\phantomsection
\addcontentsline{toc}{subsection}{B}
\subsection*{B}

\begin{itemize}
    \item \textbf{B2B:} Business-to-Business. Modello di relazione commerciale in cui le transazioni e gli scambi di beni o servizi avvengono tra aziende.
    \item \textbf{Backlog:} Elenco ordinato e prioritizzato delle attività, funzionalità e requisiti da realizzare durante il progetto.

    \item \textbf{Baseline:} Versione stabile e approvata di un insieme di documenti o artefatti, utilizzata come riferimento ufficiale per le fasi successive del progetto.

    \item \textbf{Backend:} Parte di un'applicazione che gestisce la logica applicativa, l'elaborazione dei dati e la comunicazione con servizi e sistemi esterni.

    \item \textbf{Branch:} Ramo indipendente di sviluppo all'interno di un repository Git, utilizzato per realizzare modifiche senza intervenire direttamente sul ramo principale.

    \item \textbf{Bug:} Errore o comportamento non corretto di un sistema software rispetto a quanto previsto dai requisiti o dalle specifiche.
\end{itemize}

% -------------------- C --------------------

\phantomsection
\addcontentsline{toc}{subsection}{C}
\subsection*{C}

\begin{itemize}
    \item \textbf{Capitolato:} Documento proposto da un'azienda o da un committente che descrive gli obiettivi, i requisiti e i vincoli di un progetto software.

    \item \textbf{Caso d'Uso:} Descrizione delle interazioni tra un attore e il sistema per il raggiungimento di un obiettivo specifico.

    \item \textbf{CI/CD:} Continuous Integration / Continuous Delivery (Integrazione Continua / Consegna Continua). Insieme di pratiche che automatizzano l'integrazione, la verifica e la preparazione del software per il rilascio.
    \item \textbf{Commit:} Operazione Git che registra un insieme di modifiche all'interno del repository, associandovi un messaggio descrittivo.

    \item \textbf{Committente:} Soggetto che richiede la realizzazione del prodotto software e ne definisce le esigenze e gli obiettivi.

    \item \textbf{Conventional Commits:} Convenzione per la scrittura dei messaggi di commit che utilizza una struttura standardizzata per descrivere tipologia e ambito delle modifiche.

    \item \textbf{Consuntivo:} Rilevazione delle risorse effettivamente impiegate durante un'attività, utilizzata per confrontare i risultati ottenuti con le stime iniziali.

    \item \textbf{Contingenza:} Insieme delle azioni correttive da attuare nel caso in cui un rischio identificato si manifesti.

    \item \textbf{Copertura dei Test (Code Coverage):} Metrica che misura la percentuale di codice sorgente eseguita durante l'esecuzione dei test.

\item \textbf{CRM (Customer Relationship Management):} Sistema utilizzato per gestire e organizzare le informazioni relative ai clienti e alle interazioni commerciali con essi.
\end{itemize}

% -------------------- D --------------------

\phantomsection
\addcontentsline{toc}{subsection}{D}
\subsection*{D}

\begin{itemize}
    \item \textbf{Daily Stand-up:} Breve momento di coordinamento durante il quale i membri del gruppo condividono lo stato delle attività, gli eventuali impedimenti e le necessità di supporto.

    \item \textbf{Dato non strutturato:} Informazione che non segue uno schema predefinito e che, pertanto, non può essere inserita direttamente in una struttura organizzata senza una preventiva elaborazione.
    \item \textbf{Debito Tecnico:} Insieme delle conseguenze derivanti da scelte implementative o progettuali subottimali, che possono rendere più difficile la manutenzione e l'evoluzione del software.

    \item \textbf{DoD:} Definition of Done (Definizione di Completamento). Insieme di criteri che un'attività deve soddisfare affinché possa essere considerata completata
    \item \textbf{Deployment:} Processo di distribuzione e installazione di un'applicazione in un ambiente in cui può essere eseguita e utilizzata.

    \item \textbf{Diario di Bordo:} Documento o momento di rendicontazione periodica attraverso il quale il gruppo comunica lo stato di avanzamento del progetto, le attività svolte e le eventuali difficoltà riscontrate.
\end{itemize}

% -------------------- E --------------------

\phantomsection
\addcontentsline{toc}{subsection}{E}
\subsection*{E}

\begin{itemize}
    \item \textbf{Endpoint:} Punto di accesso di un servizio o di un'API, identificato generalmente da un URL, attraverso il quale è possibile inviare richieste e ricevere risposte.

    \item \textbf{ERP:} Enterprise Resource Planning. Sistema gestionale aziendale che integra e coordina diversi processi interni, come amministrazione, contabilità, acquisti e gestione delle risorse.

    \item \textbf{Errore:} Evento che impedisce a un'operazione di concludersi correttamente o che determina un comportamento non previsto del sistema.

    \item \textbf{Euristiche:} Regole pratiche utilizzate per individuare soluzioni a problemi complessi, riducendo il tempo e le risorse necessarie per affrontarli.
\end{itemize}

% -------------------- F --------------------

\phantomsection
\addcontentsline{toc}{subsection}{F}
\subsection*{F}

\begin{itemize}
    \item \textbf{Feature:} Funzionalità o caratteristica offerta da un prodotto software.

    \item \textbf{Framework:} Insieme strutturato di librerie, strumenti e convenzioni che fornisce una base per lo sviluppo di applicazioni software.

    \item \textbf{Frontend:} Parte di un'applicazione con cui l'utente interagisce direttamente, generalmente attraverso un'interfaccia grafica.
\end{itemize}

% -------------------- G --------------------

\phantomsection
\addcontentsline{toc}{subsection}{G}
\subsection*{G}

\begin{itemize}
    \item \textbf{Git:} Sistema di controllo di versione distribuito che permette di registrare le modifiche apportate ai file e di coordinare il lavoro di più sviluppatori.

    \item \textbf{GitHub:} Piattaforma online basata su Git che consente di ospitare repository, collaborare allo sviluppo software e gestire le modifiche tramite issue e pull request.

    \item \textbf{GitHub Actions:} Servizio di automazione integrato in GitHub che permette di eseguire procedure automatiche, come compilazione, test e controlli di qualità.
\end{itemize}

% -------------------- I --------------------

\phantomsection
\addcontentsline{toc}{subsection}{I}
\subsection*{I}

\begin{itemize}
    \item \textbf{Issue:} Elemento utilizzato per descrivere, assegnare e monitorare un'attività, un problema o una richiesta all'interno di un progetto.

    \item \textbf{Iterazione:} Periodo di lavoro limitato durante il quale viene sviluppato e verificato un insieme definito di funzionalità o attività.
\end{itemize}

% -------------------- L --------------------

\phantomsection
\addcontentsline{toc}{subsection}{L}
\subsection*{L}

\begin{itemize}
    \item \textbf{LLM:} Large Language Model. Modello di intelligenza artificiale addestrato su grandi quantità di dati testuali, in grado di comprendere e generare contenuti in linguaggio naturale.

    \item \textbf{Log:} Registrazione cronologica di eventi, operazioni o messaggi generati da un sistema, utile per monitorarne il funzionamento e individuare eventuali errori.
\end{itemize}

% -------------------- M --------------------

\phantomsection
\addcontentsline{toc}{subsection}{M}
\subsection*{M}

\begin{itemize}
    \item \textbf{Manutenibilità:} Capacità di un sistema software di essere modificato, corretto o esteso con un impegno ragionevole.

    \item \textbf{Milestone:} Traguardo significativo del progetto, associato al completamento di una fase o al raggiungimento di un risultato rilevante.

    \item \textbf{Metriche:} Misure quantitative utilizzate per valutare caratteristiche del processo di sviluppo o del prodotto software.

    \item \textbf{Mockup:} Rappresentazione visiva preliminare dell'interfaccia di un'applicazione, utilizzata per mostrare la disposizione degli elementi e le principali caratteristiche grafiche.
    \item \textbf{Multimodalità:} Capacità di un sistema di elaborare e integrare informazioni provenienti da diverse modalità, come testo, immagini, documenti e audio.
\end{itemize}

% -------------------- P --------------------

\phantomsection
\addcontentsline{toc}{subsection}{P}
\subsection*{P}

\begin{itemize}
    \item \textbf{PB:} Product Baseline. Revisione formale del progetto nella quale viene fissato il prodotto software completo, insieme alla relativa documentazione e agli artefatti finali.

    \item \textbf{Piano di Progetto:} Documento che descrive l'organizzazione delle attività, la pianificazione temporale, l'assegnazione delle risorse e il monitoraggio dell'avanzamento del progetto.

    \item \textbf{Piano di Qualifica:} Documento che definisce le strategie, le attività e i criteri utilizzati per verificare e validare la qualità del prodotto e del processo.

    \item \textbf{Pull Request:} Richiesta di integrazione delle modifiche sviluppate in un branch all'interno di un altro branch. Permette di sottoporre il lavoro a revisione prima dell'integrazione.

    \item \textbf{Proponente:} Soggetto che propone il capitolato e fornisce indicazioni e chiarimenti relativi agli obiettivi e alle esigenze del progetto.

    \item \textbf{PoC:} Proof of Concept .Realizzazione sperimentale e limitata di una soluzione software, finalizzata a verificarne la fattibilità tecnica e l'adeguatezza rispetto agli obiettivi individuati.

    \item \textbf{Pull Request Review:} Attività di revisione delle modifiche contenute in una pull request, finalizzata a verificarne correttezza, qualità e conformità alle convenzioni adottate.
\end{itemize}

% -------------------- Q --------------------

\phantomsection
\addcontentsline{toc}{subsection}{Q}
\subsection*{Q}

\begin{itemize}
    \item \textbf{Qualità del Software:} Insieme delle caratteristiche che determinano la capacità di un prodotto software di soddisfare i requisiti e le aspettative degli utenti e dei committenti.
\end{itemize}

% -------------------- R --------------------

\phantomsection
\addcontentsline{toc}{subsection}{R}
\subsection*{R}

\begin{itemize}
    \item \textbf{Redattore:} Ruolo responsabile della produzione del contenuto di un documento, nel rispetto delle convenzioni e dei modelli definiti dal gruppo.

    \item \textbf{Repository:} Archivio che contiene i file di un progetto software e la relativa cronologia delle modifiche.

    \item \textbf{Requisito:} Caratteristica, funzionalità o vincolo che il sistema deve soddisfare per rispondere alle esigenze degli utenti e del committente.

    \item \textbf{Responsabile:} Ruolo incaricato di coordinare le attività del gruppo, monitorare l'avanzamento del progetto e rappresentare il gruppo nei rapporti con il committente e il proponente.

    \item \textbf{Revisione:} Attività di controllo di un documento o di un prodotto software, finalizzata a individuare errori e verificare il rispetto dei criteri stabiliti.

    \item \textbf{RTB:} Requirements and Technology Baseline. Prima revisione formale del progetto, nella quale vengono definiti e consolidati i requisiti e le tecnologie adottate, supportate da un Proof of Concept.

    \item \textbf{Rischio:} Evento incerto che, qualora si verifichi, può influenzare negativamente o positivamente il raggiungimento degli obiettivi del progetto.
\end{itemize}

% -------------------- S --------------------

\phantomsection
\addcontentsline{toc}{subsection}{S}
\subsection*{S}

\begin{itemize}
    \item \textbf{Scrum:} Framework agile per la gestione e lo sviluppo di prodotti complessi, basato su iterazioni brevi, ruoli definiti e momenti periodici di verifica e miglioramento.

    \item \textbf{Stakeholder:} Persona, gruppo o organizzazione che ha interesse nei risultati del progetto o che può influenzarne lo svolgimento.

    \item \textbf{Sprint:} Iterazione di durata definita durante la quale il gruppo si impegna a completare un insieme di attività pianificate.

    \item \textbf{Sprint Backlog:} Insieme delle attività selezionate per uno sprint e che il gruppo si impegna a portare a termine.

    \item \textbf{Sprint Planning:} Riunione nella quale il gruppo pianifica le attività da svolgere durante lo sprint successivo.

    \item \textbf{Sprint Retrospective:} Incontro al termine dello sprint durante il quale il gruppo analizza il proprio metodo di lavoro e individua possibili miglioramenti.

    \item \textbf{Sprint Review:} Incontro nel quale vengono presentati e valutati i risultati ottenuti durante lo sprint.

    \item \textbf{Stakeholder:} Soggetto che può influenzare il progetto o essere influenzato dai suoi risultati.
\end{itemize}

% -------------------- T --------------------

\phantomsection
\addcontentsline{toc}{subsection}{T}
\subsection*{T}

\begin{itemize}
    \item \textbf{Tailoring:} Processo di adattamento di uno standard o di un modello organizzativo alle caratteristiche e alle necessità specifiche del progetto.

    \item \textbf{Task:} Unità di lavoro che descrive un'attività da svolgere, generalmente associata a un responsabile, una stima temporale e una scadenza.

    \item \textbf{Template:} Struttura predefinita utilizzata come modello per la redazione di documenti o la realizzazione di artefatti, al fine di garantire uniformità e coerenza.

    \item \textbf{Test:} Attività finalizzata a verificare il comportamento di un sistema software attraverso l'esecuzione di prove e il confronto dei risultati ottenuti con quelli attesi.

    \item \textbf{Tracciabilità:} Capacità di ricostruire le relazioni tra requisiti, attività, decisioni, documenti e componenti del prodotto durante il ciclo di vita del progetto.
\end{itemize}

% -------------------- U --------------------

\phantomsection
\addcontentsline{toc}{subsection}{U}
\subsection*{U}

\begin{itemize}
    \item \textbf{UML:} Unified Modeling Language. Linguaggio standard di modellazione utilizzato per rappresentare graficamente la struttura e il comportamento di un sistema software.

    \item \textbf{Usabilità:} Grado in cui un prodotto può essere utilizzato dagli utenti per raggiungere i propri obiettivi in modo efficace, efficiente e soddisfacente.

    \item \textbf{Utente:} Soggetto che interagisce direttamente con il sistema software per usufruire delle funzionalità offerte.
\end{itemize}

% -------------------- V --------------------

\phantomsection
\addcontentsline{toc}{subsection}{V}
\subsection*{V}

\begin{itemize}
    \item \textbf{Validazione:} Processo che verifica che il prodotto realizzato soddisfi le esigenze e le aspettative del committente e degli utenti.

    \item \textbf{Verifica:} Processo che controlla la correttezza e la conformità di un prodotto o di un documento rispetto ai requisiti e alle specifiche definite.

    \item \textbf{Verificatore:} Ruolo incaricato di controllare la qualità e la conformità degli artefatti prodotti, senza coincidere con il redattore del medesimo elaborato.

    \item \textbf{Versionamento:} Gestione delle diverse versioni di un documento o di un prodotto software, in modo da mantenerne la cronologia e identificare le modifiche apportate.
\end{itemize}

% -------------------- W --------------------

\phantomsection
\addcontentsline{toc}{subsection}{W}
\subsection*{W}

\begin{itemize}
    \item \textbf{WoW:} Way of Working. Insieme di regole, strumenti, procedure e metodologie adottate dal gruppo per organizzare e svolgere il lavoro in maniera uniforme e tracciabile.
\end{itemize}